\section{Arithmetic and nested partitions}\label{sec:hierarchies}

We now turn the distance equalities in the prepared graph into an exact
identity for overlaps of distinct pairs.  The two ingredients are the
factorization of a squared planar distance and the product formula in a
number field.  Nested partitions let us express their logarithms through
nonnegative overlap lengths.

\subsection{Two coordinates at every place}

Fix one of the algebraic configurations supplied by
Proposition~\ref{prop:graph}.  Take a number field $K\subset\C$ containing
all its coordinates and $\mathrm i=\sqrt{-1}$.  For a point $x=(u,b)$ set
\[
 Z_1(x)=u+\mathrm i b,\qquad Z_2(x)=u-\mathrm i b.
\]
Both maps are injective, because $u,b$ are real in the given embedding.
Their differences remain nonzero under every field embedding of $K$.
At another complex embedding the two coordinates need not be complex
conjugates; their injectivity and the algebraic factorization below
are the properties we use.
For distinct $x,y$ the squared distance factors in $K$ as
\begin{equation}\label{eq:hier-factor}
 t(x,y)=(Z_1(y)-Z_1(x))(Z_2(y)-Z_2(x)).
\end{equation}

Let $v$ run over the places of $K$.  At a finite place above a prime $p$,
use the absolute value extending $|p|_p=p^{-1}$ and put
$\sigma_v=[K_v:\Q_p]/[K:\Q]$.  All archimedean places are complex,
since $\mathrm i\in K$; there use ordinary complex modulus and
$\sigma_v=2/[K:\Q]$.  The normalized product formula
\cite[Theorems~7.14--7.15 and Lemma~8.6]{ProductFormula} reads
\begin{equation}\label{eq:hier-product}
 \sum_v\sigma_v\log|z|_v=0\quad(z\in K\setminus\{0\}),
 \qquad \sum_{v\mid\infty}\sigma_v=1.
\end{equation}
The weights convert the norm-based finite absolute values and squared
complex moduli in that reference to the absolute values used here,
after division by $[K:\Q]$.
Here and below logarithms are natural.  For a fixed configuration only
finitely many finite places make any coordinate difference
$Z_j(x)-Z_j(y)$ a nonunit: each nonzero element of $K$ has nonzero
valuation at only finitely many prime ideals, and there are finitely
many differences.

At a place $v$, define the \emph{raw depths}
\[
 i_j(x,y)=-\log|Z_j(x)-Z_j(y)|_v\quad(x\ne y),\qquad j=1,2.
\]
The place is suppressed in this notation.  On a fiber $F_e$ with source
$x_e$ and squared distance $t_e$, Equation~\eqref{eq:hier-factor} gives
\begin{equation}\label{eq:hier-fiber}
 i_1(x_e,y)+i_2(x_e,y)=h_e,
 \qquad h_e=-\log|t_e|_v,\qquad y\in F_e.
\end{equation}
The number $h_e$ is finite and independent of the choice of $y$ in the
fiber.

\subsection{Nested grids and their error}

At a finite place, put $d_j(x,y)=i_j(x,y)$ for distinct points.  The
ultrametric inequality makes $d_j(x,y)\ge r$ an equivalence relation.
At a complex place we obtain the same nesting property from square
grids, with a controlled error in the depth.

Independently of all vertex and fiber sampling, choose independently a
uniformly random orientation and $\xi$ uniform on $[0,\log2)$.  The grid
side lengths are $l_m=e^\xi2^m$, $m\in\mathbb Z$.  Choose a shift uniform
modulo $l_0$ in each oriented axis.  To obtain the shift modulo
$l_{m+1}$ from that modulo $l_m$, for $m\ge0$, add either $0$ or $l_m$
with equal probabilities, independently at every step and in each
axis.  For $m<0$ reduce the initial shift modulo $l_m$.  Every shift is
uniform modulo its side length, and the resulting half-open square
grids are nested.  Use such a grid family for each coordinate $j$.
Randomly shifted recursive square dissections also appear in Arora's
Euclidean approximation algorithms \cite[Section~2.1]{Arora1998}.

For distinct points, let $d_j(x,y)$ be minus the logarithm of the
smallest grid side at which their images share a cell.  In both the
finite and complex cases put $d_j(x,x)=+\infty$ and
$\Delta_j(x,y)=d_j(x,y)-i_j(x,y)$ for $x\ne y$.

The three depth quantities have distinct roles:
\begin{center}
\begin{tabular}{@{}ll@{}}
\toprule
Quantity & Meaning for a distinct pair \\
\midrule
$i_j(x,y)$ & Raw logarithmic depth, $-\log|Z_j(x)-Z_j(y)|_v$ \\
$d_j(x,y)$ & Sharing depth in the nested hierarchy \\
$\Delta_j(x,y)$ & Hierarchy error, $d_j(x,y)-i_j(x,y)$ \\
\bottomrule
\end{tabular}
\end{center}

\begin{lemma}\label{lem:hier-grid}
At every complex place the depths are almost surely finite on distinct
pairs.  The distribution of $\Delta_j(x,y)$ is independent of the
distinct pair, the place, and the configuration.  Its finite mean is
an absolute constant $\kappa$, and absolute constants $C_0,b_0>0$
satisfy
\begin{equation}\label{eq:grid-error}
 \Pr\bigl(|\Delta_j(x,y)|>u\bigr)\le C_0e^{-b_0u}
 \qquad(u\ge0).
\end{equation}
At finite places $\Delta_j=0$.
\end{lemma}

\begin{proof}
Write $R>0$ for the separation of the two complex values, and
$(a,b)$ for their displacement along the oriented axes.  At side
length $l$, a uniform shift gives sharing probability
\[
 (1-|a|/l)_+(1-|b|/l)_+.
\]
Sharing is impossible if $l<R/\sqrt2$, whereas separation has
probability at most $(|a|+|b|)/l\le\sqrt2R/l$.  The separation events
decrease as the grids become coarser and their probabilities tend to
zero.  The pair therefore eventually shares a cell almost surely,
and there is a smallest sharing scale.

Write $\Delta=d_j(x,y)+\log R$.  The lower bound on a sharing side
gives $\Delta\le\log\sqrt2$.  For $u\ge0$, let $l$ be the largest grid
side no larger than $Re^u$, so $l>Re^u/2$.  On $\{\Delta<-u\}$ the
pair is separated at that side; hence
\[
 \Pr(\Delta<-u)\le2\sqrt2e^{-u}.
\]
This proves~\eqref{eq:grid-error} and a uniform first absolute moment.

For any real $u$, the event $\{\Delta\ge u\}$ is sharing at the
largest side $l\le Re^{-u}$.  The law of $l/R$ depends only on $u$,
by the uniform logarithmic phase, and the displacement direction
relative to the axes is uniform.  The sharing formula therefore
makes the entire law of $\Delta$ independent of the pair.  The
construction uses countably many random variables and cell tests,
so the depths and all these events are measurable.
\end{proof}

\subsection{Removing the giant cell}

Discard the null event on which any distinct-pair depth is not finite.
For either type of place let $\Pi_j(r)$ be the partition with
pair-sharing criterion $d_j(x,y)\ge r$, for $r\in\R$.  These partitions
are nested as $r$ increases.  Call a cell a \emph{giant} if its size
exceeds $n/2$, and denote the unique such cell, when it exists, by
$G_j(r)$.  Replace that cell by its complement and keep all other
cells.  The resulting family $\mathcal C_j(r)$ is a multiset: if the
complement is itself one other cell, it occurs twice.  Every member
$C$ has $\alpha(C)\le1/2$.
Figure~\ref{fig:giant} illustrates why multiplicities matter.
\begin{figure}[htbp]
\centering
\begin{tikzpicture}[x=0.85cm,y=0.85cm,every node/.style={font=\small}]
  \node at (2.5,1.55) {Original partition, $n=8$};
  \draw[rounded corners,fill=black!5] (0,0) rectangle (3,1);
  \draw[rounded corners,fill=blue!8] (3.4,0) rectangle (5.4,1);
  \node at (1.5,0.5) {$\{1,2,3,4,5\}$};
  \node at (4.4,0.5) {$\{6,7,8\}$};
  \node at (1.5,-0.45) {Giant $G$};
  \node at (4.4,-0.45) {$C$};
  \draw[-{Stealth},thick] (5.65,0.5) -- (6.55,0.5);
  \node at (9.3,1.55) {Transformed multiset};
  \draw[rounded corners,fill=blue!8] (6.8,0) rectangle (8.8,1);
  \draw[rounded corners,fill=blue!8] (9.4,0) rectangle (11.4,1);
  \node at (7.8,0.5) {$\{6,7,8\}$};
  \node at (10.4,0.5) {$\{6,7,8\}$};
  \node at (7.8,-0.45) {$P\setminus G$};
  \node at (10.4,-0.45) {$C$};
\end{tikzpicture}
\caption{Replacing a giant cell by its complement can duplicate an
existing cell. Both occurrences are retained in every overlap sum.}
\label{fig:giant}
\end{figure}


This replacement removes the overlap of all pairs at very coarse
scales.  Its cost is a function of each endpoint separately, which
will disappear when pair laws with the same combined marginals are
compared.  Define
\begin{align*}
 M_j&=\sup\{r:G_j(r)\text{ exists}\},\\
 S_j(x)&=\int_{\{r:G_j(r)\text{ exists}\}}
              \1_{\{x\notin G_j(r)\}}\,dr,\\
 L_j(x,y)&=\int_\R\sum_{C\in\mathcal C_j(r)}
                    \1_C(x)\1_C(y)\,dr\qquad(x\ne y).
\end{align*}

\begin{lemma}\label{lem:hier-depth}
For each realized hierarchy on $n\ge2$ points, the displayed
quantities are finite and, for distinct $x,y$,
\begin{equation}\label{eq:depth}
 d_j(x,y)=M_j+L_j(x,y)-S_j(x)-S_j(y).
\end{equation}
They are integrable over the grids.  At every finite place where all
distinct coordinate differences are units, $M_j=S_j=L_j=0$.
\end{lemma}

\begin{proof}
Suppress $j$, and let $d_-$ and $d_+$ be the minimum and maximum of
the finitely many distinct-pair depths.  Below $d_-$ the partition
has only the cell $P$; above $d_+$ it has only singletons.  A giant
at any level has a giant ancestor at every coarser level.  Thus a
giant exists on an initial ray with endpoint $M\in[d_-,d_+]$,
up to an immaterial endpoint convention.

Let $o(r)=\1_{\{d(x,y)\ge r\}}$, $g(r)=\1_{\{G(r)\text{ exists}\}}$,
and $t(r)=\sum_{C\in\mathcal C(r)}\1_C(x)\1_C(y)$.  Expanding the
giant indicator through its complement gives the pointwise identity
\[
 o(r)-g(r)=t(r)
       -\1_{\{G(r)\text{ exists},\ x\notin G(r)\}}
       -\1_{\{G(r)\text{ exists},\ y\notin G(r)\}}.
\]
The integral of the left side is $d(x,y)-M$.  Each term on the
right vanishes outside $[d_-,d_+]$ on a distinct pair, and $t(r)\le2$.
Consequently integration proves~\eqref{eq:depth}, with bounds
\[
 0\le S(x)\le d_+-d_-,\quad
 0\le L(x,y)\le2(d_+-d_-),\quad
 |M|\le\max(|d_-|,|d_+|).
\]
These bounds are grid integrable by Lemma~\ref{lem:hier-grid}, since
the configuration has finitely many pairs and finite raw depths.
At a finite place with all raw depths zero, $d_-=d_+=0$, proving the final
assertion.  Ties change neither the identity nor any integral.
\end{proof}

The empty complement of $P$ and singleton occurrences contribute
zero to every distinct-pair overlap.  We do not define $L_j(x,x)$:
singleton cells persist to arbitrarily fine scales, so its integral
would be infinite.

\subsection{The exact overlap identity}

For a nonnegative or integrable quantity depending on the place
and grids, write
\[
 \avg{H}=\sum_v\sigma_v\E_{\mathrm{grids}}H,
 \qquad
 \I(g)=\avg{\sum_{j=1}^2\int_\R
                   \sum_{C\in\mathcal C_j(r)}g\,dr}.
\]
There is no grid randomness at finite places.  The integrand $g$
may also depend on the coordinate, level, and cell occurrence.
All sums over $\mathcal C_j(r)$ retain multiplicity.  These
notations do not assert finiteness for an arbitrary $g$; finiteness
will be established for each expression used.

Set
\[
 S^*(x)=\avg{S_1(x)+S_2(x)},\qquad
 L^*(x,y)=\avg{L_1(x,y)+L_2(x,y)}\quad(x\ne y).
\]
Lemma~\ref{lem:hier-depth} and the finite set of nonunit places show
that these quantities, and $\avg{M_1+M_2}$, are finite.  For each
distinct pair, the product formula and Lemma~\ref{lem:hier-grid}
give
\[
 \avg{d_1(x,y)+d_2(x,y)}=2\kappa.
\]
Taking expectations in~\eqref{eq:depth} is therefore legitimate and
yields the promised additive identity
\begin{equation}\label{eq:additive}
 L^*(x,y)=b_*+S^*(x)+S^*(y)\quad(x\ne y),
 \qquad b_*=2\kappa-\avg{M_1+M_2}.
\end{equation}
In particular, any signed law on distinct pairs with zero total
mass and zero sum of its two marginals has zero integral against
$L^*$.  This consequence is exact; no error from sampling a
diagonal pair is involved.

The three nonnegative finite scales used below are
\begin{equation}\label{eq:hier-scales}
 W_2=\I\bigl(\pi(C^2)\bigr)=\E_\pi L^*(x,y),\qquad
 A=\E_\alpha S^*(x),\qquad W=W_2+A.
\end{equation}
Thus $W_2$ measures distinct-pair overlap after the giant is
replaced, while $A$ measures the average time spent outside the
giant.  Equation~\eqref{eq:additive} relates these scales to the
graph and fiber laws.  Every absolute grid-error bound remains
uniform in $K$, since its archimedean weights sum to one.
